Redes varejistas precisam decidir quando e quanto repor estoque em séries de demanda heterogêneas. Uma regra pode preservar a disponibilidade de um produto e gerar excesso ou falta em outro. Esta dissertação propõe o \textit{Adaptive Inventory Policy Engine} (AIPE), uma abordagem para investigar a seleção de políticas de reposição a partir de características operacionais de séries loja-produto. O \textit{Policy Selection Engine} (PSE), seu componente supervisionado, permanece em desenvolvimento e validação.

A infraestrutura integra preparação dos dados, caracterização da demanda, sinais auxiliares de previsão e comparação de políticas em simulação. A avaliação exploratória utiliza dados de uma rede varejista brasileira: um piloto com 15 lojas da Paraíba e uma expansão com 145 séries da Bahia. O portfólio reúne políticas clássicas, meta-heurísticas e aprendizado por reforço. A escolha considera custo de inventário sob exigência de serviço, com indicadores de ruptura e estabilidade dos pedidos.

Os resultados históricos indicam diferenças entre políticas e perfis, mas não demonstram generalização do seletor. Reduções de custo dependem do comparador, do serviço e da versão experimental; escolhas retrospectivas não equivalem a ganhos futuros. A contribuição é a formulação e a infraestrutura de uma avaliação contextual rastreável. Uma avaliação externa reúne 2.408 séries do M5 e evidencia diferenças de custo e serviço entre as bases, condicionadas aos respectivos protocolos. Validação temporal e treinamento do PSE são necessários para estabelecer utilidade operacional.

\vspace{2em}
\textbf{Palavras-chave}: Otimização de Inventário; Demanda Intermitente; Seleção Contextual de Políticas; Cadeia de Suprimentos.
